% 子文件：4.微分中值定理与导数应用.tex
% 主文档使用 \input{4.微分中值定理与导数应用} 引入，无document环境

\section{微分中值定理}

定理4.1（罗尔定理）\quad 若函数$f(x)$满足
\begin{enumerate}
    \item 在闭区间$[a,b]$上连续；
    \item 在开区间$(a,b)$内可导；
    \item $f(a)=f(b)$，
\end{enumerate}
则在开区间$(a,b)$内至少存在一点$\xi$，使得
\[
f'(\xi)=0.
\]

定理4.2（拉格朗日中值定理）\quad
若函数$f(x)$满足
\begin{enumerate}
    \item 在闭区间$[a,b]$上连续；
    \item 在开区间$(a,b)$内可导，
\end{enumerate}
则在开区间$(a,b)$内至少存在一点$\xi$，使得
\[
f(b)-f(a)=(b-a)f'(\xi). \tag{3}
\]

\noindent 分析（利用罗尔定理证明拉格朗日中值定理）

定理4.3（柯西中值定理）\quad 如果函数$f(x)$和$g(x)$满足
\begin{enumerate}
    \item 在闭区间$[a,b]$上连续；
    \item 在开区间$(a,b)$内可导，且$g'(x)\neq 0$，
\end{enumerate}
则在开区间$(a,b)$内至少存在一点$\xi$，使得
\[
\frac{f(b)-f(a)}{g(b)-g(a)}=\frac{f'(\xi)}{g'(\xi)}.
\]

\section{洛必达法则}
\noindent 和型未定式

\noindent 其他型未定式

\section{泰勒公式}
定理4.4（泰勒定理）\quad 设函数$f(x)$在点$x_0$处有$n$阶导数，则在$x_0$附近$f(x)$可表示为。

麦克劳林公式
\[
f(x)=f(0)+f'(0)x+\frac{f''(0)}{2!}x^2+\cdots+\frac{f^{(n)}(0)}{n!}x^n+\frac{f^{(n+1)}(\theta x)}{(n+1)!}x^{n+1}.
\]

\noindent 初等函数的麦克劳林公式
\[
\frac{1}{1-x}=1+x+x^2+\cdots+x^n+o(x^n).
\]

\section{极值与最大（小）值的求法}
\subsection{函数的极值及其求法}

\subsection{函数的最大（小）值的求法}

\section{函数的分析作图法}
\subsection{凸函数、曲线的凸向及拐点}
定理4.9\quad 设$f(x)$在区间$I$上有二阶导数，若$f''(x)\ge 0\,(\le 0)$，则$f(x)$为$I$上的下凸函数（上凸函数）。

定理4.10\quad 若

\noindent 在连续曲线$y=f(x)$上，不同凸向曲线段的分界点叫做拐点。

\subsection{曲线的渐近线}

\subsection{函数的分析作图法}

\section{曲率}
\subsection{弧微分}
若函数曲线是由极坐标方程
\[
r=r(\theta),\quad \alpha\le\theta\le\beta
\]
给出，且$r(\theta)$有连续导数，则将$x=r(\theta)\cos\theta,\ y=r(\theta)\sin\theta$代入，即得
\[
ds=\sqrt{[r(\theta)]^2+[r'(\theta)]^2}\,d\theta.
\]

\subsection{曲率}
为了说明在点$M$处曲线弯曲程度，让$M_1$沿曲线趋于$M\,(\Delta s\rightarrow 0)$，称倾角$\alpha$对弧长$s$的导数
\[
\frac{d\alpha}{ds}=\lim_{\Delta s\rightarrow0}\frac{\Delta\alpha}{\Delta s}
\]
的绝对值为曲线在点$M$处的曲率，记为$k$，则
\[
k=\left|\frac{d\alpha}{ds}\right|.
\]
就是说，曲率等于切线倾角随弧长变化的变化率的绝对值，或者说是切线方向随弧长变化的变化率的模。

在直角坐标下$\alpha=$，其微分。由式知$ds=\sqrt{1+y'^2}\,dx$就得到曲率的计算公式
\[
k=\left|\frac{y''}{\left(1+y'^2\right)^{\frac{3}{2}}}\right|.
\]

若曲线$c$上点$M$处的曲率$k\neq0$，则称$R=\dfrac{1}{k}$为曲线$c$在点$M$处的曲率半径。

在点$M$处曲线凹向的法线上取一点$D$，使$|MD|=R=\dfrac{1}{k}$，则以$D$为圆心，$R$为半径的圆叫做曲线在$M$处的